%!TeX root =  ../../nocovka.tex

\section{Commiting a git commit}

\begin{table} [H]
	\centering
	\catcode`\-=12
	\begin{tabular}[c]{|| c | c ||}
	\hline
		\multicolumn{2}{||c||}{ \textbf{Commit type cheat table}} \\
	\hline
 		 \textbf{Type} & \textbf{Repository change} \\
	\hline
		\textcolor{purple}{Feat:} & A new feature is added \\
	\hline
		\textcolor{blue}{Fixed:} & A bug is fixed \\
	\hline
		\textcolor{orange}{Docs:} & A documentation is updated \\
	\hline
		\textcolor{black}{Refactor:} & A code change that is not affecting functionality \\
	\hline
		\textcolor{brown}{Test:} & A code change in unit tests \\
	\hline
		\textcolor{gray}{Chore:} & A repository maintenance action \\
	\hline
		\textcolor{magenta}{Version x.y} & When new version is released for clearer commit history \\
	\hline
	\end{tabular}
	\caption{Types of commit headers}
	\label{table:cmtHeaders}
\end{table}